\documentclass[conference]{IEEEtran}

% --- Grundkonfiguration ---
\usepackage[utf8]{inputenc}
\usepackage[T1]{fontenc}
\usepackage[ngerman]{babel}

% --- Wissenschaftliche Pakete ---
\usepackage{cite}
\usepackage{url}

\begin{document}

\title{Zwischenabgabe: Literaturverzeichnis\\
\large Der Noisy-Neighbor-Effekt: Eine Untersuchung mikroarchitektonischer Ressourcen-Interferenzen in virtualisierten Umgebungen}

\author{%
  \IEEEauthorblockN{Maximilian Wagner}
  \IEEEauthorblockA{%
    \textit{Studiengang IT-Sicherheit} \\
    \textit{Fachbereich Informatik und Medien} \\
    \textit{Technische Hochschule Brandenburg} \\
    maximilian.wagner@th-brandenburg.de
  }
}

\maketitle

\section{Abdeckung der geforderten Bereiche}
Das vorliegende Literaturverzeichnis umfasst 12 Primärquellen, die die Grundlage für die Bearbeitung der Forschungsfrage bilden. Die Auswahl orientiert sich am Untersuchungsdesign der Arbeit und deckt die drei geforderten Bewertungskriterien ab:

\subsection{Technische Grundlagen und Methodiken}
Dieser Bereich behandelt die theoretischen Grundlagen der Virtualisierung und die Methodik für den Proof of Concept. Plauth et al. \cite{plauth2017performance} beschreiben die Unterschiede zwischen Hardware-Virtualisierung und OS-Level-Ansätzen. Felter et al. \cite{felter2015updated} und Li et al. \cite{li2017performance} liefern die Methodik zum Performance-Vergleich zwischen KVM und Containern, welche die Basis der Fallback-Strategie bildet. Danielsson et al. \cite{danielsson2019testing} stellen den experimentellen Ansatz vor, um Cache-Verdrängungseffekte auf Multi-Core-Systemen zu messen.

\subsection{Aktuelle Themen}
Der Fokus liegt auf aktuellen Hardware-Mitigierungen gegen den Noisy-Neighbor-Effekt sowie den zugehörigen regulatorischen Anforderungen. Das NIST \cite{nist2018hypervisor} stuft fehlende Plattform-Isolation in Server-Umgebungen als kritisches Risiko ein. Veitch et al. \cite{veitch2017performance} und Chintapalli et al. \cite{chintapalli2022numasfp} untersuchen die Effektivität von Mechanismen wie der Intel Cache Allocation Technology (CAT) zur Isolation von LLC-Ressourcen. Larsson et al. \cite{larsson2025esther} fassen den aktuellen Forschungsstand zu hardwaregestütztem QoS-Enforcement zusammen.

\subsection{Wissenschaftliche Publikationen}
Diese Publikationen dienen der theoretischen Fundierung der Arbeit und der Einordnung der untersuchten Schwachstellen. Koh et al. \cite{koh2007analysis} liefern grundlegende Untersuchungen zu Leistungsinterferenzen in virtuellen Umgebungen. Nikounia und Mohammadi \cite{nikounia2018hypervisor} differenzieren bei der Analyse von Performance-Einbußen zwischen Hardware-Cache-Effekten und Hypervisor-Scheduling. Ge et al. \cite{ge2018survey} ordnen Angriffe auf geteilte Caches in eine allgemeine Taxonomie ein. Gajanin et al. \cite{gajanin2025performance} untersuchen entsprechende Denial-of-Service-Risiken in aktuellen Cloud-Architekturen.

\section{Qualität und Formalia}
Die aufgeführten Dokumente umfassen Peer-Review-Publikationen (Konferenzbände, Fachjournale) und Standards offizieller Gremien (NIST). Die Quellenangaben sind formal vollständig, mit DOI-Links (sofern verfügbar) versehen und nach IEEE-Richtlinien formatiert. Die geforderte Anzahl von 8 bis 12 Quellen wird mit exakt 12 Einträgen eingehalten. Zur vollständigen Transparenz ist die zugrundeliegende BibTeX-Datenbank öffentlich im GitHub-Repository des Projekts einsehbar: \url{https://github.com/mwagner86/Hardware-Sicherheit/blob/main/paper/references.bib}.

\bibliographystyle{IEEEtran}
\bibliography{references}

\end{document}